\begin{abstract}

Video Question Answering (VQA) en videos largos plantea el desafío central de extraer información relevante y modelar dependencias de largo alcance a partir de muchos frames redundantes. El mecanismo de self-attention ofrece una solución general para el modelado de secuencias, pero tiene un costo prohibitivo cuando se aplica a una cantidad masiva de tokens espaciotemporales en videos largos. La mayoría de los métodos previos se basan en estrategias de compresión para reducir el costo computacional, como disminuir la longitud de entrada mediante sparse frame sampling o comprimir la secuencia de salida enviada al large language model (LLM) mediante space-time pooling. Sin embargo, estos enfoques ingenuos sobrerrepresentan la información redundante y a menudo omiten eventos salientes o patrones espaciotemporales de aparición rápida. En este trabajo, introducimos \model, un eficiente state-space model para manejar videos long-form. Nuestro modelo aprovecha el algoritmo selective scan para aprender a seleccionar de manera eficaz la información crítica del video de alta dimensión y transformarla en una secuencia reducida de tokens para un procesamiento eficiente por parte del LLM. Experimentos exhaustivos demuestran que \model\ alcanza precisión state-of-the-art en múltiples benchmarks de VQA long-form, incluidos PerceptionTest, NExT-QA, EgoSchema, VNBench, LongVideoBench, Video-MME y MLVU. El código y los modelos están disponibles en \url{https://sites.google.com/view/bimba-mllm}.


\end{abstract}
